\documentclass[12pt]{article}
\usepackage[utf8]{inputenc}
\usepackage[T1]{fontenc}
\usepackage{hyperref}
\title{Khumi Awa Chin Language Resources for K-12 Curriculum Development}
\author{Compiled by Claude Sonnet 5.5 (claude-sonnet-5-5) for Andr\'as Kornai}
\date{}
\begin{document}
\maketitle

\section{Introduction}
Khumi Awa Chin (ISO 639-3 cka) is a Kuki-Chin language of Chin and Arakan (Rakhine) States, Myanmar. Sources label it variously as a variety of Khumi or Mro-Khimi, so material about ``Khumi'' may not match this variety. Little was found; no school materials were located.

\section{Primary Educational Resources}
None found.

\section{Linguistic and Reference Materials}
\begin{itemize}
\item Khumi language overview: \url{https://en.wikipedia.org/wiki/Khumi_language}
\item Joshua Project language profile (Khumi): \url{https://joshuaproject.net/index.php/languages/cnk}
\end{itemize}

\section{Religious and Community Resources}
\begin{itemize}
\item Asia Harvest profile of the Khumi in Myanmar (reports a Bible in 2014, New Testament in 1959): \url{https://www.operation.asiaharvest.org/MYANMAR-Khumi.pdf}
\end{itemize}

\section{Digital and Audio Resources}
\begin{itemize}
\item Global Recordings Network, Khumi Chin: Awa (Bible lessons and audio/video): \url{https://globalrecordings.net/en/language/1150}
\item Global Recordings Network, Khumi Chin: Bangladesh (a different variety): \url{https://globalrecordings.net/en/language/5091}
\end{itemize}

\section{Curriculum Development Context}
The search found no orthography, primer or school programme for this variety. Any curriculum would require fieldwork with the community and clarification of the relation between Awa and other Khumi varieties.

\end{document}